\section{Glossaire des termes anglais et militaires}
\label{annexe-glossaire}
\begingroup
\footnotesize
\newcommand*{\Gr}[1]{#1~r\textsuperscript{o}}%
\newcommand*{\Gv}[1]{#1~v\textsuperscript{o}}%
\newcommand{\Gf}[1]{{\LJpolice #1}}%
\newcommand{\Ge}[3]{\par\noindent\hangindent=1em\hangafter=1
  \textbf{#1}\ (#2)~--- #3\par}%
\setlength{\parindent}{0pt}%
\setlength{\parskip}{1.2pt plus 0.6pt}%

\noindent Le glossaire réunit les mots et expressions anglais du carnet, les termes et sigles de l'armée britannique, ainsi que quelques termes militaires français ou latins utiles au lecteur. Il s'appuie sur le relevé de l'apparat et du commentaire historique de la présente édition. Chaque entrée donne, entre parenthèses, la ou les formes du carnet (en caractères de la transcription) et un ou deux renvois de feuillet, puis le sens que le mot a dans le contexte de l'armée britannique de 1916--1917, avec, le cas échéant, une autorité d'époque (abréviations bibliographiques p.~\pageref{annexe-bibliographie}). Les mots anglais d'usage courant sont simplement traduits. Le signe [?] marque une interprétation incertaine. Les grades abrégés, les dates et les noms de lieux relèvent du tableau des abréviations et de l'index.

\medskip
\begin{multicols}{2}
\raggedright

\Ge{afternoon ; after dinner}{\Gf{afternoon}, f.~\Gr{86} ; \Gf{After dinner}, f.~\Gv{88}}{après-midi ; après le dîner.}

\Ge{anti-aircraft battery}{f.~\Gr{91}}{batterie antiaérienne. Voir \textit{Archie}.}

\Ge{A.P.C.}{\Gf{Broche A.P.C}, f.~\Gr{70}}{sigle non résolu. \textit{Army Pay Corps} (service britannique de la solde) n'est qu'une possibilité ; il peut s'agir d'une marque commerciale.}

\Ge{Archie ; Archie Singers}{\Gf{13th Corps Archie Singers}, f.~\Gr{58} ; \Gf{XIII corps Archie}, f.~\Gr{60}}{argot militaire : canon ou obus antiaérien (Fraser et Gibbons 1925) ; les \textit{Archie Singers} sont vraisemblablement la troupe de chanteurs des batteries antiaériennes du XIII\textsuperscript{e} corps, entendue en concert au foyer Y.M.C.A.}

\Ge{Army ; Army troops ; Vth Army}{\Gf{Army troops}, f.~\Gr{47} ; \Gf{V\LJsup{th} army H.Q.}, f.~\Gv{88}}{armée ; \textit{Army troops} : troupes rattachées directement à l'armée, non à une division. La \textit{Fifth Army} (V\textsuperscript{e} armée britannique, général Gough) est le nom pris par la \textit{Reserve Army} le 30~octobre 1916 ; son quartier général est à Toutencourt.}

\Ge{A.S.C.}{\Gf{les A.S.C. men}, f.~\Gr{6} ; \Gf{A.S.C}, f.~\Gv{33}, \Gv{53}}{\textit{Army Service Corps} : service britannique des transports et du ravitaillement, équivalent du train des équipages (« Train des Équipages. Army Service Corps », Dictionary 1919). Hippolyte Combes est attaché à la 152nd Company, compagnie hippomobile du train de la 18th Division. Voir \textit{Coy}.}

\Ge{at the end of the road}{f.~\Gr{75}}{au bout de la route (emplacement des lignes de M.T.).}

\Ge{a year}{\Gf{50 £ a year}, f.~\Gv{53}}{par an : traitement annuel d'un \textit{office boy}.}

\Ge{barbed wire}{\Gf{réseaux de barbed wire}, f.~\Gr{80}}{fil de fer barbelé.}

\Ge{batman}{\Gf{To a batman for saddling up my pony}, f.~\Gr{22}}{ordonnance : soldat attaché au service personnel d'un officier (« An officer's servant », Fraser et Gibbons 1925). À distinguer du \textit{groom}.}

\Ge{battery ; B Battery ; F.A.}{\Gf{Batt\LJsup{t} F.A–}, \Gf{B Battery}, f.~\Gv{48}}{batterie d'artillerie de campagne (\textit{Field Artillery}) ; les batteries d'une même brigade d'artillerie sont désignées par des lettres. Voir R.F.A.}

\Ge{Bedfords, Bedfs}{f.~\Gv{19}, \Gr{24}}{soldats du Bedfordshire Regiment (7th Battalion, de la 54th Infantry Brigade, ravitaillée par la compagnie d'Hippolyte Combes).}

\Ge{billet, billets ; billeting officer}{\Gf{billeting officer}, f.~\Gv{1} ; \Gf{billets}, f.~\Gv{7} ; \Gf{billet}, f.~\Gr{34}}{logement chez l'habitant attribué à un militaire, et le billet de logement qui l'impose à l'habitant (Fraser et Gibbons 1925) ; français « cantonnement » (Dictionary 1919). Le \textit{billeting officer} est l'officier de cantonnement, que seconde l'interprète. À distinguer du billet de chemin de fer.}

\Ge{blacksmith}{\Gf{Blacksmith d'Herissart}, f.~\Gv{84}}{forgeron, maréchal-ferrant.}

\Ge{Bloody France !}{\Gf{Bloody France etc. !}, f.~\Gv{56}}{« Maudite France ! » : \textit{bloody}, intensif injurieux (mot du capitaine Williams).}

\Ge{Boche(s)}{f.~\Gv{8}, \Gv{53}}{Allemand(s) : terme populaire français.}

\Ge{boot polish}{\Gf{(Nugget) boot polish}, f.~\Gv{52} ; \Gf{Brown Boot Polish}, f.~\Gr{58} ; \Gf{Tin of black}, f.~\Gv{52}}{cirage (noir ou brun), en boîte ; \textit{Nugget} : marque de la Nugget Polish Co., de Londres.}

\Ge{box}{\Gf{Box of tooth powder}, f.~\Gv{19}}{boîte.}

\Ge{brainless ass}{\Gf{A brainless ass}, f.~\Gv{53}}{« âne sans cervelle » (jugement porté sur le lieutenant Curtis).}

\Ge{breakfast}{f.~\Gr{2}, \Gv{45}}{petit déjeuner, premier des repas du mess (Fernand-Laurent 1917).}

\Ge{brickfield area}{f.~\Gv{40}}{secteur des briqueteries, à Albert, où bivouaque la compagnie en octobre 1916.}

\Ge{brigade}{\Gf{refilling point de ma brigade}, f.~\Gv{33}}{brigade d'infanterie (trois ou quatre bataillons) ; la compagnie d'Hippolyte Combes ravitaille la 54th Infantry Brigade de la 18th Division.}

\Ge{cakes ; sweets}{\Gf{Sweets}, \Gf{cakes}, f.~\Gr{29}}{gâteaux ; confiseries.}

\Ge{Canada in Flanders}{\Gf{Cananda in Flanders}, f.~\Gr{22}}{titre d'ouvrage : premier volume de l'histoire officielle du Corps expéditionnaire canadien, par Sir Max Aitken (Londres, Hodder and Stoughton, 1916), acheté 1~F~25 à la cantine de la gare de Raincheval.}

\Ge{canteen ; Exp. F. Canteen}{\Gf{canteen}, f.~\Gr{22}, \Gr{59} ; \Gf{Canteen station de Raincheval}, f.~\Gr{22} ; \Gf{Exp. F. Canteen}, f.~\Gv{83}}{cantine ; le plus souvent un dépôt des \textit{Expeditionary Force Canteens} (E.F.C.), organisation officielle de l'armée britannique créée en 1915, qui vend aux troupes, près du prix coûtant, les produits britanniques (Fortescue 1928).}

\Ge{cantonnement}{\Gf{cantonnts}, f.~\Gr{54} ; \Gf{cantonnem\LJsup{t}}, f.~\Gv{75}}{logement des troupes chez l'habitant, dans les granges et les fermes ; voir \textit{billet}.}

\Ge{car ; Sunbeam Car ; car in workshop}{\Gf{cars}, f.~\Gr{49} ; \Gf{Sunbeam Car}, \Gf{car in workshop}, f.~\Gr{58}}{voiture automobile (de service) ; la \textit{Sunbeam} (Wolverhampton) du major Levy ; ailleurs, une Vauxhall (f.~\Gv{83}), voiture d'état-major courante de l'armée britannique. Voir \textit{workshops}.}

\Ge{carte d'état-major}{\Gf{Cte d Etat-major}, f.~\Gr{2}}{feuille de la carte de France au 1/80\,000 dite « d'état-major ».}

\Ge{Christmas}{\Gf{Carte : christmas}, f.~\Gr{56}}{Noël : cartes de vœux de Noël.}

\Ge{church service ; Church meeting ; Wesleyans}{\Gf{church service}, \Gf{(Wesleyans}, f.~\Gv{22} ; \Gf{Church meeting}, f.~\Gr{45}}{culte (des méthodistes wesleyens) ; réunion religieuse.}

\Ge{clerk, clerks}{f.~\Gr{47}, \Gv{48}}{commis, secrétaire de bureau.}

\Ge{coffee milk}{\Gf{2 tins of coffee milk}, f.~\Gr{24}}{lait concentré (sucré) au café, en boîte [?].}

\Ge{convoy}{f.~\Gv{25}}{convoi de voitures du train.}

\Ge{cooking ; cooks ; headcook}{\Gf{cooking}, f.~\Gv{7} ; \Gf{headcook}, f.~\Gv{3} ; \Gf{cooks}, f.~\Gv{39}}{cuisine (payée à la logeuse : rubrique de comptes) ; cuisiniers ; chef cuisinier (du mess).}

\Ge{corps}{\Gf{XIII\LJsup{e}}, f.~\Gv{52} ; \Gf{13th}, f.~\Gr{58} ; \Gf{Ve corps}, f.~\Gv{62}}{corps d'armée britannique, groupant plusieurs divisions ; les \textit{XIII Corps} et \textit{V Corps} sont nommés.}

\Ge{cottage}{f.~\Gv{1}, \Gr{2}}{maisonnette.}

\Ge{Coy, Cy ; Coy Transport}{\Gf{Coy}, f.~\Gr{5} ; \Gf{152 Coy}, f.~\Gv{47} ; \Gf{Cy}, f.~\Gv{75} ; \Gf{Coy Transport}, f.~\Gr{45}}{\textit{Company} : compagnie ; presque toujours la 152nd Company A.S.C. d'Hippolyte Combes (« ma Coy »), parfois une autre (« 150 Coy, 18\textsuperscript{th} div\textsuperscript{on} », f.~\Gv{60}). L'abréviation réglementaire est « Co. » (FSPB 1914) ; « Coy » est celle de l'usage.}

\Ge{Daily Mail}{f.~\Gr{54}}{quotidien londonien, imprimé aussi à Paris pour le continent ; le sens exact de la mention reste obscur.}

\Ge{Damn !}{f.~\Gv{7}}{juron : « Zut ! » (au départ inattendu de la compagnie).}

\Ge{darning}{\Gf{washing \& darning}, f.~\Gv{37}}{reprisage, raccommodage.}

\Ge{days}{\Gf{(2 days)}, f.~\Gv{7}}{« (pour) deux jours » (rubrique de comptes).}

\Ge{Death or Glory Boys ; 17th Lancers}{\Gf{17\LJsup{th} Lancers}, f.~\Gv{85}}{le 17\textsuperscript{e} lanciers (17th Duke of Cambridge's Own Lancers), surnommé d'après son insigne, une tête de mort avec la devise « Or Glory » (Fraser et Gibbons 1925).}

\Ge{dinner ; tea-dinner}{\Gf{dinner}, f.~\Gr{40}, \Gr{82} ; \Gf{tea-dinner}, f.~\Gr{93}}{dîner, principal repas du mess, à heure fixe (Fernand-Laurent 1917) ; \textit{tea-dinner} : repas du soir tenant lieu de thé et de dîner.}

\Ge{division}{\Gf{18\LJsup{th} div\LJsup{on}}, f.~\Gv{60} ; \Gf{63\LJsup{rd} div.}, \Gf{naval div.}, f.~\Gr{85} ; \Gf{2\LJsup{nd} australian division}, f.~\Gr{82}}{division d'infanterie (trois brigades, avec artillerie et train) : la 18th (Eastern) Division, celle d'Hippolyte Combes ; la 63rd (Royal Naval) Division, formée de marins et passée à l'armée de terre en 1916 ; la 2nd Australian Division.}

\Ge{D.O.S.T.}{\Gf{major Levy, of the D.O.S.T.}, f.~\Gr{46}}{sigle non attesté ailleurs ; sans doute la direction des approvisionnements et des transports de l'armée (\textit{Director} ou \textit{Deputy Director of Supplies and Transport}, chef de l'A.S.C. d'une armée : Puckle 1917) [?]. Voir \textit{Purchase Board}.}

\Ge{drink}{\Gf{Drink}, f.~\Gv{4}, \Gr{33}}{boisson, consommation (rubrique de comptes).}

\Ge{driver}{\Gf{Tip au driver du lorry}, f.~\Gr{34}}{conducteur.}

\Ge{dubbin}{\Gf{Tin of Wren's Dubbin}, f.~\Gv{54} ; \Gf{Wren's Dubbin}, f.~\Gr{63}}{graisse pour le cuir, qui assouplit et imperméabilise bottes et harnais (Clifton et Grimaux 1883, s.v. \textit{dubbing}) ; \textit{Wren's} : marque de Northampton.}

\Ge{dug-out(s)}{\Gf{german dug-out}, f.~\Gv{42} ; \Gf{dug-outs}, f.~\Gr{92}}{abri souterrain, cagna (Fernand-Laurent 1917).}

\Ge{eggs ; food}{\Gf{Food}, f.~\Gv{1} ; \Gf{Eggs \& food}, f.~\Gr{24}}{œufs ; nourriture, en-cas (rubrique de comptes).}

\Ge{English ; English novel}{\Gf{English}, f.~\Gr{64}, \Gv{82} ; \Gf{English novel}, f.~\Gv{79}}{les Anglais (les militaires britanniques) ; roman anglais.}

\Ge{extra}{\Gf{Food (extra)}, f.~\Gv{2}}{en supplément (rubrique de comptes).}

\Ge{farm ; wood}{\Gf{Belle-Eglise Farm}, f.~\Gr{59} ; \Gf{Belle Eglise Farm}, f.~\Gv{85} ; \Gf{acheux Wood}, f.~\Gr{84}}{ferme ; bois : toponymes sous la forme anglaise des cartes et des ordres britanniques.}

\Ge{firing line}{f.~\Gv{47}}{ligne de feu, première ligne.}

\Ge{free}{f.~\Gr{69}}{gratuit, offert.}

\Ge{French Mission ; Mission française ; M.F.}{\Gf{Mission Fr.\LJsup{se}}, f.~\Gv{18} ; \Gf{M. F\LJsup{se}}, f.~\Gv{52} ; \Gf{French Mission}, f.~\Gv{56} ; \Gf{miss. fr.}, f.~\Gv{75}}{la Mission militaire française attachée à l'armée britannique, qui affecte, administre et paie les interprètes (« solde \& indemnité », f.~\Gv{18}).}

\Ge{garret}{f.~\Gr{4}, \Gr{5}}{mansarde, galetas.}

\Ge{gift}{f.~\Gr{34}, \Gv{71}}{gratification, cadeau (au Q.M.S., à une domestique).}

\Ge{green forage}{\Gf{for green forage}, f.~\Gv{3} ; \Gf{green forage}, f.~\Gr{12} ; \Gf{green fourrage}, f.~\Gv{15}}{fourrage vert (ray-grass, luzerne, sainfoin, trèfle : FSPB 1914), fauché et payé sur place, ration d'appoint des chevaux du train.}

\Ge{groom}{f.~\Gv{2}, \Gr{6}}{palefrenier : soldat chargé du cheval d'un officier monté, en plus de l'ordonnance (Fernand-Laurent 1917) ; celui d'Hippolyte Combes est Connelly.}

\Ge{G.S. wagon}{\Gf{limber \& G.S. wagons}, f.~\Gr{37}}{\textit{General Service wagon} : chariot réglementaire à quatre roues, attelé de quatre chevaux (FSPB 1914). Voir \textit{wagon}.}

\Ge{G.V.}{\Gf{cour de G.V.}, f.~\Gv{31}}{grande vitesse : service des messageries et bagages des chemins de fer.}

\Ge{hair cut}{f.~\Gr{19}, \Gv{43}}{coupe de cheveux.}

\Ge{home ; room ; floor}{\Gf{home}, f.~\Gv{41} ; \Gf{room}, f.~\Gv{38} ; \Gf{floor}, f.~\Gv{78}}{chez soi, au logement ; chambre ; sol (de briques).}

\Ge{horse-lines}{f.~\Gr{6}, \Gv{22}}{lignes d'attache des chevaux d'une unité montée ou attelée, au cantonnement ou au bivouac.}

\Ge{H.Q., Headquarters}{\Gf{HQ\LJsup{rs}}, f.~\Gr{23} ; \Gf{Headquarters}, f.~\Gv{46} ; \Gf{H.Q.}, f.~\Gv{88}}{quartier général (d'une division, d'une armée).}

\Ge{hut ; mess-hut}{\Gf{hut}, f.~\Gr{50}, \Gr{65} ; \Gf{mess-hut}, f.~\Gv{45}}{baraque (en planches ou de fortune) ; baraque du mess. Voir Y.M.C.A.}

\Ge{infantry battalions}{f.~\Gr{38}}{bataillons d'infanterie.}

\Ge{in front}{f.~\Gr{56}}{à l'avant (de l'automobile).}

\Ge{interprète}{\Gf{int\LJsup{te}}, f.~\Gv{22} ; \Gf{interp\LJsup{te}}, f.~\Gv{49} ; \Gf{intte}, f.~\Gr{82}}{interprète militaire français détaché auprès d'une unité britannique ; voir \textit{French Mission}, R.A.T.}

\Ge{interview ; long talk ; walk ; meeting}{\Gf{meeting}, f.~\Gr{55} ; \Gf{walk}, f.~\Gr{82} ; \Gf{interview}, f.~\Gr{83} ; \Gf{long talk}, f.~\Gv{89}}{entretien ; longue conversation ; promenade ; réunion.}

\Ge{in the open}{\Gf{Cantonnés in the open}, f.~\Gv{24}}{à découvert, en plein champ.}

\Ge{khaki}{f.~\Gr{49}, \Gv{59}}{kaki, couleur des uniformes britanniques, que portent aussi les interprètes (« Casque peint khaki »).}

\Ge{killed in action ; wounded}{\Gf{killed in action}, f.~\Gr{45} ; \Gf{wounded}, f.~\Gv{81}}{tués à l'ennemi (formule des listes de pertes) ; blessés.}

\Ge{kit}{f.~\Gv{45}, \Gr{82}}{paquetage d'officier, avec la toile qui sert de sac de couchage (Fernand-Laurent 1917).}

\Ge{lecture}{\Gf{a very interesting lecture on Canada}, f.~\Gv{63}}{conférence (au foyer Y.M.C.A.).}

\Ge{Lieutenant ; Second Lieutenant}{\Gf{Lt}, f.~\Gv{1} ; \Gf{Lieut\LJsup{t}}, f.~\Gr{53} ; \Gf{Sec\LJsup{d} Lt}, f.~\Gv{89}}{lieutenant ; sous-lieutenant. Autres grades : \Gf{Cap.}, \Gf{Capt.} (\textit{Captain}, capitaine), \Gf{major}, \Gf{Col.} (colonel).}

\Ge{lift}{\Gf{eu a lift s/ 1 lorry}, f.~\Gv{39} ; \Gf{lift sur 1 water-cart}, f.~\Gv{92}}{place obtenue au passage d'un véhicule militaire ; pratique tolérée, dite \textit{lorry hopping} (Fraser et Gibbons 1925).}

\Ge{limber}{\Gf{conduit limber}, f.~\Gr{37}}{avant-train ; \textit{limbered wagon} : voiture à deux éléments articulés.}

\Ge{lorry, lorries}{\Gf{lorry}, f.~\Gr{34} ; \Gf{lorries}, f.~\Gv{33}}{camion automobile (du transport automobile de l'A.S.C.). Voir \textit{lift}, \textit{supply column}.}

\Ge{lost my way}{f.~\Gr{42}}{« perdu mon chemin ».}

\Ge{lunch ; early lunch ; Lunché}{\Gf{lunch}, f.~\Gr{2} ; \Gf{early lunch}, f.~\Gv{52} ; \Gf{Lunché}, f.~\Gv{75}}{déjeuner, collation légère vers une heure (Fernand-Laurent 1917) ; déjeuner pris tôt ; « déjeuné » (verbe francisé).}

\Ge{major}{\Gf{major}, f.~\Gr{48} ; \Gf{major Levy}, f.~\Gr{46}}{grade britannique : commandant. À distinguer du \textit{Town Major}.}

\Ge{men}{f.~\Gv{75}}{les hommes, la troupe (\textit{N.C.O's \& men}).}

\Ge{mess ; messroom ; mess-tent ; messer}{\Gf{mess}, f.~\Gr{1}, \Gv{20} ; \Gf{messroom}, f.~\Gv{7} ; \Gf{mess-tent}, f.~\Gr{39} ; \Gf{messer}, f.~\Gr{20} ; \Gf{mes off}, f.~\Gv{20}}{popote : table commune des officiers d'une unité, avec ses règles et ses rites, et le local où elle se tient (salle, baraque, tente) ; l'interprète y est admis (Fernand-Laurent 1917). \textit{Messer} : prendre ses repas au mess (verbe francisé) ; \textit{mes[s] off[icers]} : mess des officiers.}

\Ge{M.O.}{\Gf{M.O ?}, f.~\Gv{78}}{\textit{Medical Officer} : médecin militaire.}

\Ge{moving ; we are moving}{\Gf{we are moving}, f.~\Gr{24} ; \Gf{moving}, f.~\Gv{85}}{« nous partons » ; en mouvement.}

\Ge{M.P.'s}{\Gf{M.P.s}, f.~\Gr{45}}{\textit{Military Police} : police militaire britannique.}

\Ge{M.T.}{\Gf{L\LJsup{ns} de M.T.}, f.~\Gr{75}}{\textit{Mechanical} (ou \textit{Motor}) \textit{Transport} : transport automobile de l'A.S.C. ; les « lignes de M.T. » sont sans doute le parc des véhicules [?].}

\Ge{N.C.O.}{\Gf{orderly N.C.O.}, f.~\Gv{18} ; \Gf{N.C.O's}, f.~\Gr{45}}{\textit{Non-Commissioned Officer} : sous-officier ou caporal (FSPB 1914).}

\Ge{new boots, nails on}{f.~\Gr{59}}{bottes neuves, cloutage (des semelles).}

\Ge{office ; Postoffice ; Train Office ; office boy}{\Gf{office}, f.~\Gr{6}, \Gv{33} ; \Gf{Postoffice}, f.~\Gv{7} ; \Gf{Train Office}, f.~\Gv{45} ; \Gf{office boy}, f.~\Gv{53}}{bureau (de la compagnie, du R.T.O., du \textit{Purchase Board}) ; bureau de poste ; bureau du commandant du train divisionnaire ; garçon de bureau.}

\Ge{on duty}{f.~\Gv{20}}{en service commandé.}

\Ge{orderly}{\Gf{orderly N.C.O.}, f.~\Gv{18}}{de service : \textit{orderly N.C.O.}, sous-officier de service.}

\Ge{ordonnance}{\Gf{l'ordonnance}, f.~\Gv{58}}{soldat au service d'un officier (« Orderly, officer's batman », Dictionary 1919). Dans « Venu à l'ordonnance, Aveluy » (f.~\Gv{41}), le mot paraît désigner un lieu, peut-être un dépôt du service du matériel (\textit{Ordnance}) [?].}

\Ge{out riding}{f.~\Gv{60}}{sorti à cheval.}

\Ge{pamphlets}{f.~\Gv{88}}{brochures.}

\Ge{passim ; Varia}{\Gf{Varia}, f.~\Gr{1}, \Gr{12} ; \Gf{passim}, f.~\Gv{39}}{latin : divers (rubrique de comptes) ; çà et là.}

\Ge{permissionnaire}{\Gf{permiss\LJsup{re}}, f.~\Gr{64} ; \Gf{permissre}, f.~\Gr{72}}{militaire en permission.}

\Ge{planton}{f.~\Gv{26}}{soldat de service chargé de porter les ordres ou d'escorter (ici, à la gare du Bourget).}

\Ge{P.L.M.}{f.~\Gr{72}}{Compagnie des chemins de fer de Paris à Lyon et à la Méditerranée.}

\Ge{poilus}{f.~\Gr{26}}{soldats français (terme populaire).}

\Ge{pony}{f.~\Gv{2}}{poney, monture d'Hippolyte Combes.}

\Ge{Purchase Board ; Army P. Board}{\Gf{Army P. Board}, f.~\Gr{56} ; \Gf{Purchase Board}, f.~\Gr{83}}{commission des achats : service de l'A.S.C. chargé, à l'échelon de l'armée, des achats locaux de paille, de foin et de fourrage (Puckle 1917) ; Hippolyte Combes y reste attaché à la demande du major Levy.}

\Ge{Q.M.S.}{\Gf{Gift au Q.M.S Tait}, f.~\Gr{34}}{\textit{Quartermaster-Serjeant} (FSPB 1914) : sous-officier chargé des vivres, du fourrage et du matériel de l'unité ; équivalent approché : maréchal des logis fourrier (Dictionary 1919).}

\Ge{quart de place}{\Gf{1/4 place}, f.~\Gr{27}}{tarif militaire réduit sur les chemins de fer.}

\Ge{railhead}{f.~\Gv{25}, \Gv{33}}{gare de ravitaillement : point du chemin de fer où les vivres et munitions passent du train aux colonnes automobiles (FSPB 1914, « Definitions »).}

\Ge{R.A.T.}{\Gf{jeunes R.A.T.}, f.~\Gr{42}}{sigle français : réserve de l'armée territoriale, et militaire qui en relève.}

\Ge{rations}{\Gf{prendre mes rations}, f.~\Gv{7}}{vivres réglementaires.}

\Ge{refilling point}{f.~\Gv{3}, \Gr{7}}{point de ravitaillement : lieu où les camions de la \textit{supply column} transbordent les vivres dans les voitures du train divisionnaire (FSPB 1914 ; « Point de contact », Dictionary 1919).}

\Ge{R.F.}{\Gf{Hair cut, par un R.F.}, f.~\Gv{43}}{sans doute un soldat des \textit{Royal Fusiliers} [?].}

\Ge{R.F.A.}{f.~\Gr{38}, \Gv{46}}{\textit{Royal Field Artillery} : artillerie de campagne britannique.}

\Ge{R. Fusrs ; Royal Fusiliers}{\Gf{R. Fus\LJsup{rs}}, f.~\Gr{1} ; \Gf{Royal Fusiliers}, f.~\Gv{1}}{Royal Fusiliers (City of London Regiment), régiment d'infanterie ; sans doute le 11th Battalion, de la 54th Brigade.}

\Ge{R.T.O.}{f.~\Gv{25}, \Gv{33}}{\textit{Railway Transport Officer} : commissaire militaire de gare (Dictionary 1919) ; c'est à son bureau qu'Hippolyte Combes dépose ses colis et apprend le déplacement de sa division.}

\Ge{R.2}{\Gf{intte Valet, R.2}, f.~\Gr{82}}{sigle non résolu.}

\Ge{saddle soap}{f.~\Gr{18}}{savon de sellerie, pour l'entretien des cuirs.}

\Ge{shift : shifté, shifterons}{\Gf{shifté}, f.~\Gv{7}, \Gv{33} ; \Gf{shifterons}, f.~\Gr{37}}{changer de cantonnement (\textit{to shift}) : verbe francisé.}

\Ge{shirt ; vest ; socks ; handk.}{\Gf{1 shirt. 1 vest / 1 p. socks. 1 handk\LJsup{ch}}, f.~\Gv{17}}{chemise ; gilet de corps ; (une paire de) chaussettes ; mouchoir (\textit{handkerchief}).}

\Ge{staff}{f.~\Gv{88}}{membre du personnel (du bureau de la Mission).}

\Ge{stars ; sunset ; sunshine}{\Gf{sunset}, \Gf{sunshine}, f.~\Gr{12} ; \Gf{stars}, f.~\Gr{24}}{coucher de soleil ; soleil, éclaircies ; étoiles.}

\Ge{straw \& hay}{f.~\Gr{44}}{paille et foin, objets des achats locaux.}

\Ge{stuck}{f.~\Gv{90}}{embourbé, immobilisé.}

\Ge{supply column}{\Gf{lorries (supply column)}, f.~\Gv{33}}{colonne de ravitaillement automobile, qui porte les vivres du \textit{railhead} au \textit{refilling point} (Puckle 1917).}

\Ge{Sussex}{f.~\Gv{11}}{Royal Sussex Regiment, régiment d'infanterie.}

\Ge{Taubes}{\Gf{Taubes}, f.~\Gr{26} ; \Gf{taubes}, f.~\Gr{73}}{allemand \textit{Taube}, « colombe » : nom du monoplan allemand du début de la guerre, appliqué à tout avion allemand.}

\Ge{tea}{f.~\Gr{25}, \Gr{80}}{le thé, repas de fin d'après-midi.}

\Ge{the fields}{\Gf{to the fields for green forage}, f.~\Gv{3}}{aux champs.}

\Ge{tin(s) of}{\Gf{2 tins of}, f.~\Gr{24}, \Gv{83}}{boîte(s) en fer-blanc.}

\Ge{tip}{f.~\Gv{2}, \Gr{20}}{pourboire.}

\Ge{tommy, tommies}{\Gf{tommies}, f.~\Gv{14} ; \Gf{tommy}, f.~\Gv{83}}{soldat britannique (de \textit{Tommy Atkins}, nom fictif du modèle de livret du soldat ; Fraser et Gibbons 1925).}

\Ge{tooth powder ; tooth ache}{\Gf{Calvert Tooth Powder}, f.~\Gr{59} ; \Gf{Calvert's Tooth Powder}, f.~\Gv{83} ; \Gf{tooth ache}, f.~\Gr{80}}{poudre dentifrice (Calvert, de Manchester) ; mal de dents.}

\Ge{Town Major}{\Gf{Town-Major}, f.~\Gv{45} ; \Gf{Town Major's office}, f.~\Gr{81}}{major de la place : officier chargé, dans une localité, du cantonnement de toutes les unités de passage (Fernand-Laurent 1917) ; il attribue et retire les \textit{billets}.}

\Ge{trench mortar school}{\Gf{trench mortar}, \Gf{school}, f.~\Gr{20}}{école de mortiers de tranchée.}

\Ge{very}{\Gf{very belle}, f.~\Gr{12} ; \Gf{a very interesting lecture}, f.~\Gv{63}}{très.}

\Ge{wagon, waggon}{\Gf{waggon}, f.~\Gv{3} ; \Gf{wagon}, f.~\Gv{18} ; \Gf{wagons}, f.~\Gr{37}}{chariot hippomobile du train (voir G.S., \textit{limber}), non le wagon de chemin de fer.}

\Ge{washing}{f.~\Gv{1}, \Gv{17}}{blanchissage (rubrique de comptes).}

\Ge{water-cart}{f.~\Gv{92}}{voiture-citerne à eau.}

\Ge{W.C.}{f.~\Gr{27}}{\textit{water-closet} : toilettes.}

\Ge{whole day off}{\Gf{a whole day, off}, f.~\Gr{61}}{une journée entière de repos.}

\Ge{workshops}{\Gf{car in workshop}, f.~\Gr{58} ; \Gf{workshops}, f.~\Gv{84}, \Gv{90}}{ateliers de réparation du transport automobile (Puckle 1917).}

\Ge{Y.M.C.A. ; Y.M.C.A. hut}{\Gf{Y.M.C.A.}, f.~\Gr{50} ; \Gf{Y.M.C.A-hut}, f.~\Gr{58} ; \Gf{Y.M.C.A. Hut}, f.~\Gr{60}}{\textit{Young Men's Christian Association} : association chrétienne qui tient auprès des troupes des foyers (\textit{huts}) au triangle rouge, où l'on achète, lit, écrit et assiste à des concerts et des conférences (Yapp 1918).}

\Ge{£}{\Gf{50 £ a year}, f.~\Gv{53}}{livre sterling (25,22~F au pair : March 1925).}

\end{multicols}
\endgroup
